\chapter{Evaluation}
\label{ch:evaluation}

\section{Evaluation strategy and its justification}
\label{sec:eval-strategy}

% "The aim" collided with sec:aims, which states a broader project aim.  This is the
% evaluation question, not the project aim, and the two must not share a name.
The project set out to measure what adaptor functionality \emph{costs}, so the evaluation
was built around one requirement: every measured difference had to have a single
attributable cause. Four decisions follow from it.

Correctness comes before cost. Performance figures for an incorrect implementation are
worthless, and no formal proof was in scope, so correctness rests principally on
differential evidence: a pinned known-answer value reproduced byte for byte by a second,
independently written implementation. Internal self-consistency is not enough, and one
error described below left every primitive passing in isolation while the assembled whole
was wrong.

The primary comparison is controlled. LAS is measured against its own simplified base at
identical algorithm, parameters and primitives, so the paired difference estimates the
cost of the adaptor layer and of nothing else. Measuring against optimised
CRYSTALS-Dilithium instead would confound that cost with parameter, packing and
hint-vector differences; that contrast is therefore kept only as context, and the
classical ECDSA adaptor only as a functionality-matched price of post-quantum.

% NOT "two tiers": sec:benchmethod declares THREE boundaries -- core, packed, and the
% hybrid one the classical library exposes.  The warrant for "misleads quietly" is local
% ("twice here"), never a general claim about benchmarking.
Every measurement declares its boundary, because an undeclared boundary misleads quietly,
as it did twice here before the tiers were fixed. For the same reason each run gates a
directly counted rejection rate against its closed-form prediction rather than inferring
acceptance from a timing ratio.

Metrics follow the venue. The application study reports execution time and communication
bytes, counting off-chain messages, because the evaluated UTXO ledger carries no gas
metric and because off-chain messaging is where the post-quantum cost lands.

\section{Achievement of the objectives}
\label{sec:eval-objectives}

\begin{table}[htbp]
  \centering
  % ⚠ This table is referenced from 01-introduction.tex ("All five objectives were met"),
  % so tab:objectives must not be deleted: doing so leaves a dangling \cref in Chapter 1,
  % which is also why the caption no longer repeats that verdict.  Its pointers ARE the
  % point of it -- a marker without the repository needs to know where each claim is
  % evidenced -- so the row text is what compresses, never the references.
  \caption[Objectives against evidence]{Each objective of \cref{sec:aims} against the
           evidence for it, which spans the introduction, \cref{ch:results} and the
           appendix.}
  \label{tab:objectives}
  \begin{tabular}{@{}p{0.24\textwidth}p{0.68\textwidth}@{}}
    \toprule
    \textbf{Objective} & \textbf{Evidence} \\
    \midrule
    1. Literature and scheme selection
      & Adaptor signature chosen for its direct relevance to atomic swaps; LAS
        \cite{esgin2020post} for the shortest path from a standardised lattice signature
        (\cref{sec:relatedwork}). \\[3pt]
    2. Additive implementation
      & Lattice core reused byte-for-byte with no upstream function modified, audited by
        a two-branch diff; adaptor layer, wire encoding and application are new code
        (\cref{tab:reuse}). \\[3pt]
    3. Validation
      & Functional, serialization, tamper and known-answer tests pass
        (\cref{sec:res-correctness}); an independent Rust implementation reproduces the
        wire format and the pinned value byte-for-byte (\cref{sec:res-rust}). \\[3pt]
    4. Fair benchmarking
      & Adaptor overhead isolated per operation at two declared tiers
        (\cref{tab:overhead-l3}), against two further baselines, with the rejection rate
        gated against theory and the relative conclusions reproduced by the Rust port
        (\cref{tab:rust}). \\[3pt]
    5. Atomic-swap demonstration
      & End-to-end two-party, two-chain swap settling both legs and asserting
        unspendability before adaptation (\cref{sec:res-app}), extended to a measured
        three-configuration study on a UTXO ledger (\cref{sec:res-stage2}). \\
    \bottomrule
  \end{tabular}
\end{table}

% SCOPE: only configurations 2 and 3 share one public matrix and byte-identical inputs
% (backend.rs "configurations()": one LasParams, hence one A, key material from one master
% seed).  Configuration 1 is secp256k1 and has no lattice matrix at all, so an unscoped
% sentence claims a control the study does not have -- sec:res-stage2 says "only the 2->3
% step is controlled" on its own page.
% PAGINATION: pulls one line from the next page back onto this one, so the section
% heading that follows is not left with only two lines beneath it.  This paragraph
% begins on that page, which is where \enlargethispage takes effect.
% ⚠ Pagination-dependent -- re-check on rebuild.
\enlargethispage{\baselineskip}
Fairness was not simply asserted. The standalone evaluation holds algorithm, parameters
and primitives fixed; in the application study the step from Groth16 to LaZer shares one
public matrix and byte-identical inputs, verified at run time, while the classical
configuration rests on a different curve and is a whole-stack contrast rather than a
controlled one. The figures are then defended from two independent directions: a
closed-form gate on the rejection rate, and re-measurement under a different compiler and
an independent statistical harness. That second check is also where the claims stop.
Absolute per-operation times do not reproduce across the two implementations and are not
offered as a comparison between them; what reproduces, and what the conclusions rest on,
is the relative cost of the adaptor layer.

% tab:onchain's own caption says its rows are --gas-report CHARGES, "not transaction
% totals and not client measurements", so setting the figure straight against a
% per-transaction cap conflates the two accounting boundaries that caption separates.
The demonstration was built in the UTXO setting assumed in \cite{esgin2020post}, then
measured against a gas-metered venue for contrast. There a \emph{complete}, validated
native LAS verifier was charged $\approx$\gasLasVerifiedM{} million gas for its claim
path, beyond the EIP-7825 per-transaction cap, though that figure is a
% "bound the IMPLEMENTATION rather than the scheme" is an equivalence claim, and the
% equivalence holds only under the registration invariant no contract checks (tab:onchain),
% which is why that reference stays: without it the caveat has no warrant in this chapter.
% The single instance belongs here too: two verifier stages branch on coefficient values.
\texttt{--gas-report} charge rather than a transaction receipt. The cap turned out to
bound the \emph{implementation} rather than the scheme: a restructured verifier was mined
by a real client inside it, at the target parameter set, for the message length and single
signature instance measured, and only under the registration invariant that
\cref{tab:onchain} records as unchecked. The venue decision rests on cost, not on
impossibility.

\section{Implementation challenges}
\label{sec:eval-challenges}

% REWRITTEN 2026-09-06 as continuous synthesis.  The previous version framed the section
% around "the six problems of tab:challenges" and three recurring themes; both were
% scaffolding rather than content, and the cross-reference made a body section depend on an
% appendix table (now deleted).  Each challenge is stated as problem -> why it resisted
% ordinary debugging -> fix, and every reference points at a main chapter.
%
% ⚠ SCOPES THAT MUST NOT DRIFT, each corrected during drafting:
%  - NOT "all four adaptor operations still succeed": app:challenges named THREE (PreSign,
%    PreVerify, Extract).  "the adaptor operations can still succeed" avoids the count.
%  - NOT "no local test observes it": an unbounded universal.  Scope it to the tests run.
%  - NOT "restores the adaptor properties" for ML-DSA: that experiment is a FUNCTIONAL
%    demonstration only and the security of committing to HighBits(w+Y) is unanalysed, so
%    the claim is scoped to the tested contract.
%  - NOT "the known-answer value exposed divergences": its purpose is cross-checking, and
%    nothing on record says it caught a specific divergence.  Say what it made possible.
%  - ORDER: the Solidity error was found by comparing the recomputed challenge digest
%    against the C reference; correction, then stagewise validation, then the gas-cap
%    restructuring followed.  Stagewise validation must not be named as what exposed it.
%  - The patched rule's SECURITY is unanalysed (a caveat that never lapses).  It is carried
%    here by "experimental" plus sec:res-txstruct, whose own prose states it in full.
\textsc{PreSign} must reject at $\gamma-\kappa-1$, one tighter than \cref{eq:signbound}.
At the looser bound the adaptor operations can still succeed while the adapted signature
later exceeds the norm ordinary \textsc{Verify} accepts; local adaptor tests did not expose
it, and accounting for the witness norm identified the tighter bound. A related assumption,
that Dilithium's hint and bit-splitting optimisations had to be disabled, also proved too
strong: carrying the commitment path on $w+Y$ satisfies the tested adaptor contract with
them enabled (\cref{sec:res-mldsa}).

Implementing the scheme independently in Rust required byte-identical agreement, including
deterministic sampling and serialization; the pinned known-answer value made that agreement
directly testable (\cref{sec:rustport}). Benchmarking raised a separate problem, because
rejection sampling makes sign-class timings right-skewed: acceptance inferred from a timing
ratio was biased, and the two timed loops initially performed unequal work. Attempts are
therefore counted directly, checked against theory, and measured at declared boundaries
(\cref{sec:benchmethod}).

The Groth16 circuit initially cloned growing linear combinations while constructing dense
rows, making the work quadratic in row width; building each row once through the
lower-level constraint interface reduced it to linear cost. In Solidity, an
already-NTT-domain public matrix was transformed again, producing a verifier that disagreed
with the C reference; comparing the recomputed challenge digest against C exposed the
error, the corrected port was validated in stages, and the verifier was then restructured
to fit the per-transaction gas cap (\cref{sec:res-evm}). Bitcoin imposed different
constraints: Script has no native lattice-signature verifier, and its \btcChunkWidth\,B
element limit prevents LAS signatures and public keys from being supplied as single
elements. The integration therefore chunks these objects and reassembles them inside an
experimental patched consensus rule (\cref{sec:res-txstruct}).

\section{Limitations}
\label{sec:eval-limitations}

% "three", not "two": the ZKNox dependency below is the third.  It was previously carried
% only in the deleted app:challenges caption, and it scopes a MEASURED gas figure, so it
% belongs in the body rather than an appendix (2026-09-06).
The threats of \cref{sec:res-validity} bound what the results \emph{mean}, and three further
limitations belong to the artefact itself, not to the measurement. The EVM verifiers reuse
lattice primitives from a library its authors describe as not production-ready, so the gas
figures characterise a numerically-correct verifier rather than a deployable one. The
hint-free
construction yields larger keys and signatures than optimised Dilithium: cryptographic
optimisation traded for a transparent and testable artefact. \Cref{sec:res-mldsa}
quantifies that price and shows it recoverable in principle, but only for the signature
and not for the statement.

% NOT "not a stack anyone should deploy": a normative absolute the evidence does not
% carry.  What is supported is that it is not fully post-quantum.
Configuration~2 gives up something different. Groth16's soundness rests on a pairing
assumption Shor breaks, so a post-quantum signature does not protect a swap whose fairness
rests on a classical proof. That configuration is the intermediate point which makes the
controlled proof-system comparison possible, and is reported as that and nothing more.
